% year: תשפ"ג
% type: מבחן
% semester: א
% moed: ב
% number: 6ב
% points: 10
% categories: improper-integrals
% source: exams/מבחנים/תשפג/מועד ב תשפג.docx

\begin{question}
(10 נק') בדקו התכנסות של האינטגרל $\displaystyle\int_1^{\infty}\frac{\sqrt{x^5}\,dx}{2x^2+3x+40}$.
\end{question}

\begin{hint}
בדקו כיצד מתנהגת הפונקציה כאשר $x\to\infty$ והשוו ל-$x^{p}$ מתאים (מבחן ההשוואה הגבולי).
\end{hint}

\begin{solution}
הפונקציה $g(x)=\frac{x^{5/2}}{2x^2+3x+40}$ רציפה וחיובית ב-$[1,\infty)$. נשווה ל-$h(x)=x^{1/2}$:
\[ \lim_{x\to\infty}\frac{g(x)}{h(x)}=\lim_{x\to\infty}\frac{x^2}{2x^2+3x+40}=\frac12\in(0,\infty). \]
לפי \textbf{מבחן ההשוואה הגבולי}, $\int_1^\infty g$ ו-$\int_1^\infty x^{1/2}dx$ מתכנסים או מתבדרים יחד. $\int_1^\infty x^{p}dx$ מתכנס רק עבור $p<-1$, ו-$p=\frac12$, לכן $\int_1^\infty x^{1/2}dx=\infty$.
(אפשר גם לשים לב שאפילו $g(x)\to\infty$, ולכן בוודאי שהאינטגרל מתבדר.)

\textbf{תשובה:} האינטגרל \textbf{מתבדר} (שווה $+\infty$).
\end{solution}
